\documentclass[11pt,reqno]{amsart}

\usepackage[
  letterpaper,
  left=1.2in,
  right=1.2in,
  top=1.1in,
  bottom=1.4in
]{geometry}
\usepackage[final]{microtype}
\usepackage[hidelinks]{hyperref}
\hypersetup{
  pdftitle={Teaching Statement},
  pdfauthor={Wenrui Huang}
}

\setlength{\parindent}{1.5em}
\setlength{\parskip}{0pt}
\linespread{1.05}
\emergencystretch=1em
\raggedbottom

\markboth{Teaching Statement}{Wenrui Huang}

\begin{document}

\begin{center}
{\LARGE Teaching Statement}\bigskip

{\large Wenrui Huang}
\end{center}
\vskip 22pt

Students often come to me after they have made some progress on a problem but
cannot see what should come next. At that point, I try not to give them the
next line immediately. We may go back to a definition or draw the object we
are discussing. I remember a conversation from one of my terms as a teaching
assistant for Topology at the University of Science and Technology of China
(USTC). A student kept treating the quotient space as though it were a subset
of the space we started with. Repeating the definition did not help. For about
ten minutes, we drew two disjoint circles, paired corresponding points, and
wrote beside the picture what a point of the quotient actually was: an
equivalence class. After a while, the student stopped choosing a representative
as if it were the point and began talking about the whole class. I had been
moving between representatives and equivalence classes so automatically that
I had forgotten this was a step at all.

My teaching has ranged from Topology, Differential Geometry, and Riemannian
Geometry at USTC to Calculus II and III at Brown, including a calculus course
for physics and engineering students. I used to think these subjects demanded
two quite different approaches. After teaching both kinds of course, I am less
sure. Calculus students also need to know why a method applies, and in topology
a drawing is sometimes what makes a definition usable. At Brown, this often
means connecting a formula to the object it describes. When we reverse the
order of integration over the region between
$y=x^2$ and $y=x$, we first draw the region and compare a vertical slice with a
horizontal one. Only then do we write the two iterated integrals. I ask
students to use the picture to tell me where each slice begins and ends; the
bounds become a record of something they can see instead of a pattern to
memorize. In the course for physics and engineering students, I also ask what
quantity an answer represents and whether its sign or units make sense.

I have probably learned the most about how students get stuck during office
hours. When a student brings me a calculation or proof that has gone off
course, I ask them to mark the last line they genuinely understand. We sit
side by side and work forward from there. Sometimes the missing step is local;
sometimes we have to go back and rebuild an idea from several weeks earlier. I
have learned not to take over the pencil too quickly. A student can follow a
complete correction, nod at every line, and still be unable to restart the
problem alone the next day. I ask one question and then wait. The silence can
feel long, but I want the student to decide what the next line should do.
Before the student leaves, I ask them to talk through the approach once more
in their own order. Sometimes this takes most of the appointment.

\newpage

Generative AI has made this problem harder to ignore. A student can now arrive
with a polished solution without understanding why it works, and the solution
may still be wrong. I do not think the answer is to pretend these tools are
absent. One exercise I would like to use is to give students an AI-generated
solution and ask where they first stop trusting it and what they would check.
I can also see a use for AI when a student wants another example or a different
explanation. But the student still has to decide whether the mathematics holds
together.

I use a similar approach when I read proofs. Students arrive in topology and
geometry with very different amounts of practice, and many are still learning
what a complete argument has to make explicit. When a proof breaks down, I
point to the first unsupported sentence and ask what would justify it. I keep
the comment specific: this sentence needs a reason. When the same gap appears
in several solutions, I take it as a sign that my own explanation needs
another pass. I have also had proofs that looked transparent in my notes
become much less clear when I tried to explain them without skipping steps.
Student questions often uncover an assumption I had stopped noticing or a
change of viewpoint that I had compressed into one line. Preparing an
explanation forces me to unpack what I might otherwise call standard. I now
notice the same issue in my research writing.

My research is in partial differential equations, nonlinear dispersive
equations, and kinetic equations. In an advanced PDE course, I would want to
spend time on why we are trying a particular estimate, not only on the
technical steps once the estimate has been chosen. Much of my own research
proceeds by breaking a large problem into smaller questions and keeping track
of what remains unresolved. I would be glad to teach calculus, differential
equations, real analysis, geometry, and courses in partial differential
equations. Not every student I teach will go on in mathematics. I hope they
leave my class better able to say where they are stuck and work out a next
step when a familiar method fails.

\end{document}
